\section{Regular-cover gluings as congruence problems}
\label{sec:gluings}

\subsection{The geometric model and why regularity matters}

Fix $G=C_m$ or $G=D_m$. We write dihedral elements as
$(\epsilon,t)=T^tR^\epsilon$, where $\epsilon\in\{0,1\}$,
$t\in\Z/m\Z$, and
\begin{equation}
 (\epsilon,t)(\eta,u)
   =(\epsilon\mathbin\oplus\eta,\;t+(-1)^\epsilon u).
 \label{eq:dihedral-product}
\end{equation}
For the cyclic group only parity zero is allowed. The order of $D_m$ is
$2m$, including the small degenerate parameters $m=1,2$.

Each block is a connected regular cover of a specified compact connected
orientable surface with nonempty boundary. Its fibre over a chosen basepoint
is identified with the set $G$, and monodromy acts by left multiplication.
The based peripheral system includes fixed paths from that basepoint to the
boundary circles. These paths are part of the coordinate framing.
The images of the supplied canonical free generators must generate all of
$G$; the implementation checks this by gcds and the presence of a reflection
in the dihedral case. This condition makes the block connected and its
monodromy action regular.

Any automorphism of one such block covering the identity on its fixed base
surface acts on the fibre by right multiplication. Indeed, a fibre map $f$
commuting with every left multiplication satisfies
$f(g)=gf(1)$ for all $g\in G$. Write it $g\mapsto gh$.
This observation is the reason one group element per block suffices.
It does not hold in this form for an arbitrary nonregular cover.

An oriented attachment edge $e:u\to v$ joins distinct available boundary
ports, possibly of the same block. Every port is used at most once. Choose
compatible reference gluings reversing the base boundary orientations. In
the fixed fibre coordinates, the two peripheral monodromies must be exact
inverses; this is checked. A phase $a_e\in G$ then glues the entire preimage
of one boundary circle to the entire preimage of the other by
$g\mapsto ga_e$ at the reference fibre. This is well-defined because right
multiplication commutes with the left peripheral action.

Thus attachments are allowed to permute the lifted boundary circles in a
specified algebraic fashion. They are not arbitrary independently chosen
permutations of individual lifts. The assembly graph is labelled and fixed,
as are the local monodromies, base surfaces, and boundary ports.
In particular, the natural dihedral action on $m$ points is generally not
the regular action on $2m$ points used here; substituting it would invalidate
the one-group-element description of all local transports.

\begin{remark}[Local and global regularity]
Every block cover is regular. The cover of the assembled base surface need
not be regular: compatible deck maps are constrained by the holonomy around
assembly cycles. The word ``regular'' in this section always describes the
local blocks and their available right multiplications.
\end{remark}

\subsection{Eliminating all tree phases}

Consider two phase assignments $a_e$ and $a'_e$ on identical local block data.
An isomorphism over the fixed labelled base decomposition consists of right
multiplications $h_v$ such that
\begin{equation}
 a_eh_v=h_ua'_e\qquad(e:u\to v).
 \label{eq:gauge}
\end{equation}
This equation follows by applying the map before or after the attachment to
an arbitrary fibre element. It is therefore both necessary and sufficient
for the local maps to glue to a global covering isomorphism.

Choose a spanning forest. For a component root $o$, let $A_v$ be the ordered
product of original phases on the tree path from $o$ to $v$, using inverses
when an edge is traversed backwards. Define $A'_v$ similarly.

\begin{theorem}[Forest reduction of gluing equivalence]
\label{thm:forest}
Every solution of the tree-edge equations is determined uniquely by one
root element $h=h_o$ and has the form
\begin{equation}
 h_v=A_v^{-1}hA'_v.\label{eq:forest}
\end{equation}
For a remaining edge $e:u\to v$, define
\[
 c_e=A_u a_e A_v^{-1},\qquad
 c'_e=A'_u a'_e(A'_v)^{-1}.
\]
The remaining conditions are exactly
\begin{equation}
 c_eh=hc'_e.\label{eq:simultaneous-conjugacy}
\end{equation}
\end{theorem}
\begin{proof}
For a tree edge traversed from $u$ to $v$, solve \eqref{eq:gauge} for $h_v$ and proceed
along the unique tree path. Reversing an edge gives the inverse formula, so
induction establishes \eqref{eq:forest} and uniqueness. Substitute \eqref{eq:forest} into \eqref{eq:gauge}, multiply
on the left by $A_u$ and on the right by $(A'_v)^{-1}$, and obtain \eqref{eq:simultaneous-conjugacy}.
The same manipulations in reverse prove sufficiency on all edges.
\end{proof}

All products here involve fixed-size group records and integers reduced
modulo $m$. The algorithm never unfolds a fibre of $m$ or $2m$ sheets.
There is also no need to try $|G|^v$ possible collections of local maps.

\subsection{At most two scalar congruence systems}

For $D_m$, write $h=(\epsilon,t)$ and fix $\epsilon$ temporarily.
If $c=(\eta,b)$ and $c'=(\eta',b')$, equation \eqref{eq:simultaneous-conjugacy} first requires
$\eta=\eta'$. If $\eta=0$, it requires
\begin{equation}
 b=(-1)^\epsilon b'\pmod m,
 \label{eq:rotation-condition}
\end{equation}
with no constraint on $t$. If $\eta=1$, it requires
\begin{equation}
 2t=b-(-1)^\epsilon b'\pmod m.
 \label{eq:reflection-condition}
\end{equation}
Let $d_2=\gcd(2,m)$. The latter is inconsistent unless its right side is
divisible by $d_2$. Otherwise it is a single congruence for $t$ modulo
$m/d_2$, obtained by dividing by $d_2$ and inverting $2/d_2$.
Only the two choices of $\epsilon$ are tried. The cyclic case has only
$\epsilon=0$ and all cycle equations are equality tests of rotation shifts.

Markings can be included without adding a branching factor per mark.
An ordered fibre-point mark $x\mapsto y$ at block $v$ requires
$xh_v=y$, and hence
\begin{equation}
 h=A_vx^{-1}y(A'_v)^{-1}.
 \label{eq:point-mark}
\end{equation}
A lifted boundary circle is an orbit $H x$ of the peripheral subgroup
$H=\langle p\rangle$. Requiring this named circle to map to $H y$ is
equivalent to
\begin{equation}
 xA_v^{-1}hA'_v y^{-1}\in H.
 \label{eq:boundary-mark}
\end{equation}
This condition preserves the specified circle without choosing a particular
point on it.

If $p=(0,r)$ is a rotation, then
$H=\{(0,jd):j\in\Z\}$, where $d=\gcd(m,r)$.
For fixed root parity, \eqref{eq:boundary-mark} first tests a parity and then is one congruence
for $t$ modulo $d$. If $p=(1,r)$ is a reflection, then
$H=\{(0,0),(1,r)\}$. The parity of the left side selects which of those
two elements is possible, leaving one congruence modulo $m$. These
derivations also handle the identity peripheral element: each boundary
circle then contains just one point of the reference fibre.

\begin{lemma}[Congruence intersection with an explicit obstruction]
\label{lem:crt}
The two conditions $t\equiv r\pmod d$ and $t\equiv s\pmod e$ are
compatible if and only if
$\gcd(d,e)\mid(s-r)$. If compatible, their intersection is one residue
class modulo $\operatorname{lcm}(d,e)$, obtainable by the extended Euclidean
algorithm. If incompatible, the failed divisibility is a certificate of
inconsistency.
\end{lemma}
\begin{proof}
Substitute $t=r+dz$. The equation $dz\equiv s-r\pmod e$ is solvable
exactly when the stated gcd divides its right side. Divide by the gcd to
obtain an invertible coefficient. The solutions differ by
$e/\gcd(d,e)$ in $z$, so by the lcm in $t$.
\end{proof}

\begin{theorem}[Complete marked gluing solver]
\label{thm:gluing}
For two checked assemblies in the model above and any explicit collection
of fibre-point and boundary-lift marks, all covering isomorphisms over the
fixed labelled base decomposition can be encoded by at most two records
per connected assembly graph. Each record has the form
\begin{equation}
 h_o=(\epsilon,r+jD),\qquad 0\leq j<m/D,
 \label{eq:root-family}
\end{equation}
where $D\mid m$. Every local transport follows from \eqref{eq:forest}. Empty branches
have independently checkable contradiction records. The algorithm has
polynomial bit complexity in the full listed input and $\log(m+1)$.
\end{theorem}
\begin{proof}
The forest theorem is exact. For each possible root parity, \eqref{eq:rotation-condition}--\eqref{eq:boundary-mark}
reduce every remaining requirement to a failed parity/equality test or a
single congruence modulo a divisor of $m$. Repeated use of
\cref{lem:crt} yields either a contradiction or precisely one residue class
modulo a divisor $D$ of $m$. Its representatives in $\Z/m\Z$ are \eqref{eq:root-family}.
There are no other solutions, since every original equation or marking was
included. Conversely, substitution proves that every reported family member
gives an isomorphism. A forest pass and a linear number of gcd, modular
inverse, and group operations use integers of $O(\log(m+1))$ bits, together
with the explicit graph and monodromy indices. Polynomial bit complexity
follows. Components are retained as a product of solution families rather
than enumerated jointly.
\end{proof}

The implementation provides two replay routes. One checks a single collection
of local transports directly against \eqref{eq:gauge} and the original marks, without
using a spanning forest or CRT. The other checks the reported complete
solution families and their contradiction records. It verifies the
congruences and their combined modulus, rather than rerunning the producer's
CRT search. These routes separate the correctness of a particular witness
from the stronger assertion that no isomorphisms were omitted.

For a more explicit bound, let $N$ count the listed vertices, edges, free
generators and marks, and let $B$ include the largest input-integer bit
length and $\log(N+2)$. After input normalization, every root residue and
modulus has $O(\log(m+1))$ bits. A conservative schoolbook bound for a
connected graph is $O(NB^3)$ bit operations, with $O(NB)$ bits of
representation. Constraint work is grouped by graph component, so the
solution-family computation has the same bound for a disconnected graph.
The implementation additionally reports the product of all component
isomorphism counts; explicitly multiplying those binary counts costs at
most another $O(N^2B^2)$ operations. No map family is enumerated.

\subsection{Canonicalization and exact residual state counts}

For an unmarked connected assembly, forest normalization makes the tree
phases identities. The remaining $\beta$ phases form a tuple in $G^\beta$;
the residual root change simultaneously conjugates this tuple. This gives
both a canonical cache representation and a count of inequivalent framed
assignments.

For $C_m$, conjugation is trivial. For $D_m$, if all tuple entries are
rotations, compare the tuple of shifts with its simultaneous negation and
choose the lexicographically smaller one. If there is a reflection, choose
the first. Its shift $a$ can be normalized to $\delta=0$ when $m$ is odd,
and to $\delta=a\bmod2$ when $m$ is even. For root parity $\epsilon$,
solve
\[
 2t=a-(-1)^\epsilon\delta\pmod m.
\]
There are at most two parity candidates for the resulting normalized tuple.
When $m$ is even, the two possible solutions differing by $m/2$ have the
same conjugation action because the half-turn is central. Taking the smaller
tuple therefore gives a polynomial-time canonical representative for this
unmarked, fixed-base equivalence relation.

\begin{theorem}[Orbit counts after all block deck changes]
\label{thm:orbits}
On a fixed connected assembly graph of cycle rank $\beta$, with compatible
local blocks and no additional marks, the number of phase assignments up to
block deck changes is
\begin{equation}
 N_{C_m}(\beta)=m^\beta.
 \label{eq:cyclic-count}
\end{equation}
For the regular dihedral model it is
\begin{equation}
 N_{D_m}(\beta)=
 \begin{cases}
 \dfrac{(2m)^\beta+(m-1)m^\beta+m2^\beta}{2m},&m\text{ odd},\\[5pt]
 \dfrac{2(2m)^\beta+(m-2)m^\beta+m4^\beta}{2m},&m\text{ even}.
 \end{cases}
 \label{eq:dihedral-count}
\end{equation}
\end{theorem}
\begin{proof}
Forest normalization identifies the desired classes with simultaneous
conjugacy orbits on $G^\beta$. This proves \eqref{eq:cyclic-count}. For a finite group,
double-count pairs consisting of a tuple and a group element fixing it:
the number of orbits is $|G|^{-1}\sum_{g\in G}|C_G(g)|^\beta$.
For odd $m$, the identity has centralizer size $2m$, the other $m-1$
rotations have centralizer size $m$, and each of the $m$ reflections has
centralizer size 2. For even $m$, the identity and half-turn have centralizer
size $2m$, the other $m-2$ rotations have size $m$, and reflections have
size 4. Substitution proves \eqref{eq:dihedral-count}, including $m=1,2$ and $\beta=0$.
\end{proof}

\begin{corollary}[When the assembled cover is regular]
\label{cor:global-regularity}
For a connected assembly graph, without imposing additional marks, the
deck group of the cover of the assembled base is isomorphic to the
simultaneous centralizer
\[
 \bigcap_{e\notin\text{forest}} C_G(c_e).
\]
The assembled cover is regular if and only if every chord holonomy $c_e$
belongs to the centre of $G$.
\end{corollary}
\begin{proof}
Set the source and target phase assignments equal in
\cref{thm:forest}. Root elements defining deck maps are exactly those
commuting with every $c_e$. Taking inverses accounts for the right-action
composition convention and identifies the deck group with this
centralizer. The assembled cover is connected and has degree $|G|$.
The deck action on a fibre is free, so it is transitive exactly when the
centralizer has order $|G|$, or equivalently is all of $G$. This is precisely
the stated centrality condition.
\end{proof}

Consequently every cyclic assembly in this model is globally regular.
For $D_m$ with $m\geq3$, the permitted central holonomies are just the
identity for odd $m$, and the identity and half-turn for even $m$.
The degenerate groups $D_1$ and $D_2$ are abelian, so all their assemblies
are globally regular. Marks can reduce the group of mark-preserving
automorphisms without changing whether the underlying cover is regular;
the criterion refers to the unmarked covering map.

\begin{corollary}[A precise enumeration obstruction]
\label{cor:gluing-obstruction}
An algorithm that explicitly lists one record for each cyclic framed gluing
class requires at least $m^\beta$ records. In particular, $\beta=1$ and
$m=2^b$ already give $2^b$ inequivalent classes, although each individual
phase and each equivalence query use only polynomially many bits in $b$.
\end{corollary}

This obstruction has a carefully limited meaning. It is not a lower bound
for unknot recognition, nor even for deciding an arbitrary property of the
assemblies. In fact, the connected underlying orientable surfaces have the
same Euler characteristic and boundary counts for fixed block and port data;
they need not be distinguished by unmarked surface homeomorphism. What grows
is the number of \emph{framed covering} classes. A hierarchy algorithm may
be able to discard some of that data, but it must prove that doing so
preserves the relevant attachments and continuations.

For an explicit graph, $\beta\leq e$ and the binary orbit count has
$O(e\log(m+1))$ bits, so exact exponentiation is polynomial in the listed
input. If a standalone caller supplies an arbitrary binary integer $\beta$,
the output itself can have exponentially many bits in $\log\beta$.
The orbit-count API must then be interpreted as output-sensitive.

\subsection{Surface topology without sheet expansion}

For each connected component of the assembly graph, its assembled covering
surface is connected because every local block is connected and every
attachment is nonempty. If the base blocks are $S_v$, then
\[
 \chi(\widetilde S)=|G|\sum_v\chi(S_v).
\]
Gluing along complete boundary circles changes no Euler characteristic.
An unused boundary port with peripheral element $p$ contributes
$|G|/\operatorname{ord}(p)$ boundary components. Thus
\begin{equation}
 b(\widetilde S)=\sum_{\text{unused ports }j}
     \frac{|G|}{\operatorname{ord}(p_j)},\qquad
 g(\widetilde S)=\frac{2-b(\widetilde S)-\chi(\widetilde S)}2.
 \label{eq:surface-topology}
\end{equation}
Right deck maps preserve the lifted orientation, and the base attachments
reverse boundary orientations, so the assembled surface is orientable.
All quantities in \eqref{eq:surface-topology} are evaluated in binary. They help audit the input
and output but do not replace the finer gluing equivalence conditions.

For a concrete nontrivial attachment, take two pair-of-pants bases with
local monodromies $T,R$ in $D_m$ and glue their reflection ports. Each
connected regular block cover has genus zero and $2m+2$ boundary circles.
The full preimages of the reflection ports consist of $m$ circles, so
the assembly has genus $m-1$ and $2m+4$ remaining boundary circles.
All its phases are equivalent because its assembly graph is a tree.
Thus even an assembly of exponentially large genus can have a small
exact comparison problem when $m$ is given in binary. The tests check this
nontrivial peripheral gluing directly against expanded sheet permutations.
